\chapter{Discussion}
\label{chap:discussion}

\section{The central finding: exploration and decision separate}

The results support a simple but consequential conclusion: producing a correct answer somewhere in a multi-agent trace is not the same as delivering it. Independent candidates increase the chance that useful content appears, especially for Ensemble and Parallel Roles, yet their oracle-final gaps show that a final language-model decision frequently fails to exploit that content. Debate provides the clearest audited instance. Its revised candidate set can contain a correctly endorsed answer, but the final picker loses that state at a weighted rate of 5.45\% and does not produce an audited beneficial wrong-to-correct transition. Sequential shows the opposite final-boundary pattern: destructive aggregation is rare and the aggregator sometimes repairs a wrong last-worker state.

This does not establish that Sequential is globally better than Debate. Sequential's upstream discovery is weaker, particularly on MuSiQue, and an aggregator cannot recover evidence that never enters the chain. Debate frequently achieves stronger final accuracy and high semantic discovery. The result is instead a decomposition of capability. Debate appears better at candidate generation but less reliable at final selection; Sequential has a weaker source state but a more constructive final aggregator. Architecture choice should therefore depend on which component is limiting performance.

The distinction also clarifies why final accuracy can mask design problems. A Debate system may finish correctly even after a harmful revision if the other candidate remains correct. A Sequential aggregator may finish correctly despite a wrong last worker by synthesising earlier findings. Conversely, a system may have high discovery and low final accuracy. Transition analysis exposes these compensating pathways rather than assigning all credit or blame to the architecture as a whole.

\section{Why Debate can lose correct answers}

Debate assumes that exposure to an alternative argument helps an agent revise errors and helps a picker choose the better position. The machine screen suggests both beneficial and harmful revisions, but the targeted audit shows that the judge substantially overstates some of these transitions. The final human audit nevertheless provides strong evidence that the picker is not reliably corrective in the sampled population.

Several mechanisms are plausible. First, revised outputs can converge linguistically even when evidence remains asymmetric. A confident wrong answer may supply a clean narrative that is easier for the picker to recognise than a correct but qualified answer. Second, the picker receives model-generated summaries rather than a protected evidence graph. It may evaluate rhetorical coherence instead of entailment from the oracle context. Third, agents share the same model family and prompts, so their errors are not independent. Debate between correlated agents can amplify a shared shortcut rather than supply genuine epistemic diversity. Finally, the picker is a further generative call with its own exposure to truncation, extraction, and formatting errors.

These mechanisms should be treated as interpretations rather than directly identified causes. The audit labels correctness transitions, not latent persuasion. However, they lead to concrete design tests. A picker could be required to cite the exact evidence span supporting its choice. Candidate identities and confidence could be hidden to reduce stylistic bias. The system could preserve all answer-bearing claims in a structured table before revision. A deterministic evidence-based selector could also be compared against the generative picker. Each intervention targets selection without changing candidate generation, allowing the oracle-final gap to be tested directly.

\section{Why Sequential aggregation can recover}

Sequential's beneficial final transitions indicate that a wrong last-worker endorsement does not necessarily mean the accumulated trace lacks the answer. The aggregator sees the plan and worker history and may combine a correct earlier finding with the final subproblem result. In that sense it operates as a synthesiser rather than a simple selector. This explains why Sequential final semantic accuracy can exceed the designated candidate-oracle measure on some 2Wiki cells.

The architecture still has serious upstream weaknesses. Workers are serially dependent, later prompts contain growing context, and a final worker may focus on its assigned relation rather than restating the full answer. The machine screen's worker-chain losses and relatively low discovery rate suggest that relevant evidence can disappear before aggregation. The appropriate intervention is therefore different from Debate. Sequential systems need explicit state persistence: a structured ledger of established entities and relations, provenance for each intermediate fact, and constraints that prevent later workers from replacing validated findings without contrary evidence.

The beneficial aggregator result also cautions against judging intermediate workers by final-answer criteria alone. A worker assigned one hop may be correct and useful without endorsing the complete gold answer. Discovery and Endorsement labels help, but role-aware utility labels would be an even better future measure for partial subtask outputs.

\section{More computation is not necessarily more reasoning}

Doubling the nominal budget produces the expected resource increase but not an accuracy increase. Aggregate discovery and final semantic accuracy both decline, with more correct-to-wrong than wrong-to-correct paired changes. The pooled McNemar results are exploratory because they do not adjust for repeated appearances of the same question across conditions, but the descriptive direction is consistent across the aggregate counts and no architecture-level final-accuracy comparison shows a conventional 0.05-level improvement. This pattern is consistent with the argument that matched or greater token budgets do not guarantee an MAS advantage \citep{tran2026single}. It also shows where additional work can be unproductive in this dataset: longer trajectories are associated with lower discovery and greater post-discovery loss.

Three explanations are compatible with the data. First, additional tokens may be spent repeating context or generating stylistic elaboration rather than new evidence. Second, a longer allowance creates more opportunities for revision drift: a model can reason past a correct concise answer into a speculative alternative. Third, realised token use is endogenous. Hard questions and unstable states may naturally consume more tokens, so the association between longer output and error can arise without tokens causing the error. The paired cap experiment provides stronger evidence that the larger allowance did not help aggregate performance, but it still does not identify which of these mechanisms dominates.

The practical implication is to allocate compute conditionally. A system should stop when evidence and answer endorsement are stable, and spend additional budget only when candidate disagreement or missing hops justify it. Confidence alone may be poorly calibrated, so stopping should use observable conditions: agreement on an evidence-backed entity, complete coverage of plan steps, and absence of unresolved contradictions. This is more defensible than treating a maximum token count as a target.

\section{Implications for evaluation}

The project demonstrates why an evaluator needs both scale and auditability. The Mistral judge is useful: it processes tens of thousands of stages, highlights candidate failure points, and guides a focused human sample. Its broad stage agreement around 80\% is sufficient for exploratory summaries. It is not sufficient for declaring exact transition prevalence, especially when errors concentrate in the source state and when the target claim depends on rare flips.

The targeted calibration failure is itself a methodological result. A general validation set can be misleading when the downstream analysis conditions on hard cases selected by the same judge. This resembles spectrum bias: evaluator performance changes with the composition of the evaluated population. Future LLM-as-judge studies should validate on the strata used for substantive inference, not only on an overall random set. For transition analysis, validation should be pair-level because the derived class depends on both labels.

Separating Discovery from Endorsement also improves construct clarity. A single ``semantic correctness'' label can count an answer mention as success even when the stage rejects it. The audit found 14 such Discovery-correct, Endorsement-wrong rows. This is a small fraction of all annotations but directly concentrated in the mechanism of interest. An evaluator optimised for ordinary final answers may therefore fail exactly where debate and aggregation need to be studied.

\section{Implications for system design}

The results suggest four design principles.

\begin{enumerate}[leftmargin=*]
  \item \textbf{Protect discovered evidence.} Once a stage supplies an evidence-grounded candidate, store the claim and supporting span in a persistent structure. Later agents may challenge it, but should not silently erase it.
  \item \textbf{Separate generation from selection.} A selector should be evaluated independently on frozen candidate sets. Candidate coverage and picker accuracy answer different questions and should be optimised separately.
  \item \textbf{Use architecture-specific supervision.} Debate revision requires training on evidence-sensitive disagreement, while Sequential requires state tracking and synthesis. Generic instruction tuning may improve one component while leaving another unchanged.
  \item \textbf{Allocate compute adaptively.} Additional agents or tokens should be triggered by missing evidence or unresolved disagreement, not used uniformly. Stable answers with clear evidence should terminate early.
\end{enumerate}

The project's completed \texttt{v3b} checkpoint is largely aligned to single-agent behaviour. A natural extension is role-conditioned SFT using successful planner, worker, and aggregator traces. LoRA makes this feasible at modest training cost \citep{hu2021lora}, while Optima offers a related generate-filter-train perspective for multi-agent optimisation \citep{chen2024optima}. Such training must be evaluated with held-out traces and the same preservation metrics. Otherwise, an apparent gain in Sequential could be purchased by degrading the SAS or by teaching outputs that exploit the evaluator.

\section{Relationship to prior findings}

The findings neither wholly confirm nor reject prior debate research. \citet{du2023debate} and \citet{liang2023mad} demonstrate settings in which multi-agent interaction improves outcomes. The current 2Wiki results also include MAS gains, showing that the mechanism can work. The matched-budget MuSiQue results and the observed selection losses qualify the generality of that advantage. Model scale, task difficulty, candidate independence, and selector strength all matter.

The results align more directly with \citet{tran2026single}: once thinking budgets are controlled, a coherent single-agent trajectory remains a strong baseline and extra calls are not automatically productive. Building on that paper's theoretical, context-degradation, and qualitative failure accounts, this dissertation contributes an operational decomposition for the present setting: some MAS traces make answer-equivalent content available and then lose it at a measured final boundary. Communication quality and aggregation fidelity are therefore central variables in a compute-matched comparison, without implying that this is the only explanation for the observed performance gaps.

The work also extends self-consistency research \citep{wang2023selfconsistency}. Multiple samples can increase coverage, but the method of aggregation determines whether that coverage translates into accuracy. A learned generative aggregator is more flexible than majority voting but also less transparent. Reporting oracle accuracy, diversity, and regret alongside final accuracy makes this trade-off visible.

\section{Recommended next experiment}

The highest-value next experiment is a controlled selector replacement. Freeze the existing Debate revised candidates and compare three final decisions: the current generative picker, deterministic evidence matching against the oracle context, and a human selection. Because candidate generation remains identical, differences estimate selector quality without confounding from new proposals. The primary outcome should be human-audited \ctow{} and \wtoc{} rates, supplemented by final normalised accuracy and latency.

For Sequential, the corresponding experiment is a structured evidence ledger. Workers write entities, relations, and citations into immutable fields while retaining a free-text explanation. The final aggregator receives both the ledger and the transcript. A factorial ablation can test transcript only, ledger only, and both. This directly evaluates whether state persistence reduces worker-chain loss and whether the aggregator's beneficial transitions are maintained.

These experiments are preferable to immediately scaling the model or increasing the number of agents. They target the failure boundaries already identified and can reveal whether the architecture, rather than base capability, is the limiting factor.
